\documentclass[11pt,a4paper]{article}
\usepackage[T2A]{fontenc}
\usepackage[utf8]{inputenc}
\usepackage[russian,english]{babel}
\usepackage{amsmath,amssymb,amsthm,mathtools}
\usepackage{setspace}
\usepackage{enumitem}
\usepackage[hidelinks]{hyperref}
\onehalfspacing
\hyphenpenalty=10000\exhyphenpenalty=10000\emergencystretch=3em
\newcommand{\head}[1]{\par\bigskip\noindent\textbf{\underline{#1}}\quad}
\newcommand{\rudate}{14 сентября 2026 года}
\newcommand{\sheettitle}[3]{% title, subtitle, homework-flag (empty or "Homework")
  \begin{center}
  {\huge\bfseries #1\par}\vspace{1.2em}
  {\Large #2\par}\vspace{0.8em}
  \if\relax\detokenize{#3}\relax\else{\Large #3\par}\vspace{0.8em}\fi
  {\foreignlanguage{russian}{\rudate}\par}
  \end{center}\vspace{0.5em}}
\newcommand{\src}[1]{\unskip\nobreak\hfil\penalty50\hskip1em\hbox{}\nobreak\hfill(#1){\parfillskip=0pt\par}}
\newcommand{\fl}[1]{\left\lfloor #1\right\rfloor}
\DeclareMathOperator{\lcm}{lcm}
\DeclareMathOperator{\ord}{ord}
\begin{document}
\sheettitle{Number theory -- 2}{Orders and primitive roots}{}

\head{Theoretical minimum.} The multiplicative group of residues is the engine of elementary
number theory: almost every divisibility statement about $a^n\pm1$ is a statement about the order
of $a$.

\head{Definitions.} For $\gcd(a,n)=1$ the \emph{order} $\ord_n(a)$ is the least $d>0$ with
$a^d\equiv1\pmod n$. A \emph{primitive root} modulo $n$ is an element of order $\varphi(n)$, that is
a generator of $(\mathbb Z/n\mathbb Z)^*$.

\head{Theorem 1 (the basic property).} $a^m\equiv1\pmod n$ if and only if $\ord_n(a)\mid m$; in
particular the order divides $\varphi(n)$. \emph{Proof:} divide $m$ by the order with remainder.
\emph{Corollary:} if $p\mid a^n-1$ then $\ord_p(a)\mid n$ and $\ord_p(a)\mid p-1$, so every prime
divisor of $a^n-1$ either divides some $a^d-1$ with $d\mid n$, $d<n$, or is congruent to $1$ modulo
$n$. This one corollary settles most olympiad questions about the prime divisors of $2^n-1$ and of
Fermat numbers.

\head{Theorem 2 (primitive roots modulo a prime).} For every prime $p$ the group
$(\mathbb Z/p\mathbb Z)^*$ is cyclic. \emph{Proof:} for $d\mid p-1$ the congruence $x^d=1$ has at
most $d$ roots, since a polynomial of degree $d$ over the field $\mathbb F_p$ has at most $d$ roots;
hence the number of elements of order exactly $d$ is either $0$ or $\varphi(d)$; summing
$\sum_{d\mid p-1}\varphi(d)=p-1$ forces it to be $\varphi(d)$ for every $d$, and in particular
$\varphi(p-1)>0$ elements of order $p-1$ exist.

\head{Theorem 3 (the general structure).} $(\mathbb Z/n\mathbb Z)^*$ is cyclic exactly for
$n=1,2,4,p^k,2p^k$ with $p$ an odd prime; for $n=2^k$, $k\geqslant3$, it is
$\mathbb Z/2\times\mathbb Z/2^{k-2}$. In general the Chinese remainder theorem gives
$(\mathbb Z/n\mathbb Z)^*\cong\prod(\mathbb Z/p_i^{k_i}\mathbb Z)^*$, so the exponent of the group
is lcm of the individual orders --- the Carmichael function $\lambda(n)$, which can be much smaller
than $\varphi(n)$.

\head{Fact 1 (index calculus).} Fixing a primitive root $g$ modulo $p$ turns multiplication into
addition: every residue is $g^k$ with $k$ well defined modulo $p-1$, and $x^m\equiv a$ becomes a
linear congruence in the index. In particular $a$ is an $m$-th power modulo $p$ if and only if
$a^{(p-1)/\gcd(m,p-1)}\equiv1$, and the $m$-th powers form a subgroup of index $\gcd(m,p-1)$.

\head{Fact 2 (Zsygmondy).} For $a>b\geqslant1$ coprime and $n\geqslant2$ the number $a^n-b^n$ has
a prime divisor dividing no $a^k-b^k$ with $k<n$, apart from the exceptions $(a,b,n)=(2,1,6)$ and
$a+b$ a power of two with $n=2$. The statement may be quoted; it converts many existence questions
into a one-line check.

\medskip\noindent\rule{\linewidth}{0.4pt}

\begin{enumerate}[leftmargin=2.2em,itemsep=0.6em]
\item Prove Theorem 1 and its corollary, and deduce: every prime divisor of $2^p-1$ with $p$ prime
is congruent to $1$ modulo $2p$; for $k\geqslant2$ every prime divisor of the Fermat number
$2^{2^k}+1$ is congruent to $1$ modulo $2^{k+2}$; and there are infinitely many primes congruent to
$1$ modulo $n$ for every fixed $n$.

\item Prove that $(\mathbb Z/p\mathbb Z)^*$ is cyclic, counting the elements of each order, and
determine for which $n$ a primitive root modulo $n$ exists. Then show that $2$ is a primitive root
modulo $3^k$ for every $k$, and find the order of $2$ modulo $2^k-1$.

\item Let $p>3$ be a prime. A sequence $a_1,\dots,a_{p-1}$ of integers is \emph{wonky} if the terms
are distinct modulo $p$, the partial sums $a_1a_2,\ a_1a_2+a_2a_3,\ \dots,\
a_1a_2+a_2a_3+\dots+a_{p-1}a_1$ are distinct modulo $p$, and
$a_ia_{i+2}\not\equiv a_{i+1}^2\pmod p$ for all $i$, indices taken cyclically. Determine for which
$p$ a wonky sequence exists. \src{ICMC}

\item Let $n$ be an integer and $p$ a prime greater than $2$. Show that
\[
(p-1)^n\,n!\;\Bigm|\;(p^n-1)(p^n-p)(p^n-p^2)\cdots(p^n-p^{n-1}),
\]
and explain the right-hand side as the order of $\mathrm{GL}_n(\mathbb F_p)$ (Abstract
algebra -- 4). \src{CIIM 2014}

\item Find all primes $p>5$ for which there exist an integer $a$ and an integer $r$ with
$1\leqslant r\leqslant p-1$ such that the sequence $1,a,a^2,\dots,a^{p-5}$ can be rearranged into a
sequence $b_0,b_1,\dots,b_{p-5}$ with $b_n-b_{n-1}-r$ divisible by $p$ for
$1\leqslant n\leqslant p-5$. \src{Putnam 2024}
\end{enumerate}
\end{document}
